\section{Equality and multipliers on the original data}
\label{va:section}

Throughout this section the data are the original data of
Theorem~\ref{thm:rigidity}.  In particular, the densities have the geometric
normalization of Definition~\ref{def:data}.  Set
\begin{equation}\label{va:constants}
 r_h=\sqrt{A/(4\pi)}=m+\sqrt{m^2-Q^2},\qquad
 c=\frac{1-Q^2/r_h^2}{r_h}>0,\qquad
 b^0=\frac E m,\quad b^i=-\frac{P_i}{m}.
\end{equation}
For varying data write
\(\mathcal L=b^0E_{\mathrm{new}}+\sum_i b^iP_{i,\mathrm{new}}\), with
the coefficients in Equation~\eqref{va:constants} fixed.  Then
\(m'=\mathcal L'\) at the original data.  The two charge fluxes will be
fixed.  On the branch \(r>Q\), differentiating
\((r+Q^2/r)/2\), with \(a=4\pi r^2\), gives
\begin{equation}\label{va:mass-area-derivative}
 16\pi\frac{d}{da}\left(\frac{\sqrt{a/(4\pi)}+
 Q^2/\sqrt{a/(4\pi)}}2\right)=c\quad\text{at }a=A.
\end{equation}

We first allow a more general obstacle.  Let \(B\) be a smooth full cut
whose exterior is connected and contains the given end.  On each component
fix \(\varepsilon\in\{1,-1\}\), and assume
\(\theta_\varepsilon=H+\varepsilon\operatorname{tr}_B K=0\).
Either all of \(S\) is a component of \(B\), or \(B\) is disjoint from
\(S\); all other components are in \(\operatorname{int}\Omega\).
When \(B=S\) take \(\varepsilon=1\).  A compact filling is used only to
describe sets containing the obstacle, as in
Proposition~\ref{ct:perimeter-equivalence}.  No electromagnetic field is extended
into that filling.

The minimum enclosing area for \(B\) is at least \(A\), since every
enclosure of \(B\) encloses \(S\).  The mixed-sign numerical comparison,
Theorem~\ref{thm:numerical}, gives the reverse inequality at the
unchanged mass and charges.  Consequently its infimum is \(A\), and the
data exterior to \(B\) also attain equality.  We carry out the following
argument for either obstacle.

\subsection{The area derivative and a positive constraint multiplier}

Let \(\mathcal M_B\) be the family of perimeter-minimizing filled sets,
identified up to null sets and represented by their regular frontiers.
Proposition~\ref{ct:strict-compactness} confines this family, and the minimizers
for all sufficiently small smooth variations considered below, to a common
compact set.  Write \(T_M(x)\) for the unoriented tangent plane to its
frontier.  For a metric variation \(h\), define
\begin{equation}\label{va:sheet-area}
 a'_M(h)=\frac12\int_{\partial M}\operatorname{tr}_{T_M}h\,dA_g.
\end{equation}
The topology on \(\mathcal M_B\) is convergence in measure.  On this
family it implies strict perimeter convergence and convergence of the
tangent-plane area measures; in particular Equation~\eqref{va:sheet-area}
is a continuous function of \(M\).

\begin{lemma}\label{va:area-envelope}
For a smooth path with initial metric derivative \(h\), the right
derivative of its enclosing-area infimum is
\begin{equation}\label{va:area-derivative}
 a'_+(0)=\min_{M\in\mathcal M_B}a'_M(h).
\end{equation}
\end{lemma}
\begin{proof}
This is the envelope formula of Proposition~\ref{ct:area-derivative}; the
following two inequalities specify how it is used here.  On the common
compact truncation, the area density has the expansion
\[
 dA_{g_t}|_T=\left(1+\tfrac t2\operatorname{tr}_T h+O(t^2)\right)dA_g|_T,
\]
where the remainder is uniform in the point and the two-plane \(T\).
Every fixed \(M\in\mathcal M_B\) is a competitor at positive \(t\), giving
the upper derivative bound.  Conversely choose minimizing sets \(M_t\)
at positive \(t\).  Metric comparison and the same expansion show
\(P_g(M_t)\to A\).  A subsequence converges in measure to a member
\(M_0\in\mathcal M_B\), and its perimeter convergence is strict.
Continuity of Equation~\eqref{va:sheet-area} gives
\[
 a(t)-A\ge t\,a'_{M_t}(h)+O(t^2),
 \qquad \liminf_{t\downarrow0}\frac{a(t)-A}{t}\ge a'_{M_0}(h).
\]
The uniform remainder is bounded by a constant times the uniformly
bounded perimeter.  Taking a subsequence realizing the lower limit proves
Equation~\eqref{va:area-derivative}.
\end{proof}

Put \(\alpha_1=*_{g}\mathcal E^\flat\) and
\(\alpha_2=*_{g}\mathcal B^\flat\).  Use variations of
\((g,K,\alpha_1,\alpha_2)\) that are smooth and compactly supported,
with support allowed to meet \(B\); both form variations are closed.
The unit ball is transported using the positive metric square root, so
for \(g'=h\) a transported vector has derivative
\(-\tfrac12h^\sharp v\).  Let
\begin{equation}\label{va:active-rays}
 \mathcal A_B=\{(x,v):x\in\overline{\Omega_B},\ |v|_g\le1,
                  \ \mu_m+J_m(v)=0\},\qquad
 \mathfrak C'=2(\mu_m+J_m(v))'.
\end{equation}
Here \(\Omega_B\) is the exterior to the chosen obstacle.

We add four directions supported away from \(B\).  The metric direction
equals \(r^{-1}\delta\) sufficiently far out.  The three tensor directions
have Euclidean trace reversals
\begin{equation}\label{va:momentum-prototype}
 p_{ij}=r^{-2}\bigl(a_i n_j^\delta+a_jn_i^\delta
                         -\delta_{ij}a\cdot n^\delta\bigr),
 \qquad a\in\mathbb R^3.
\end{equation}
Thus their actual \(K\)-variations are
\(p-\tfrac12(\operatorname{tr}_\delta p)\delta\).
The flat linearized scalar curvature of \(r^{-1}\delta\) vanishes,
and \(\partial^jp_{ij}=0\).  Their flux variations are, respectively,
\(E'=1/2\) and \(P'_i=a_i/2\).  These four directions span the four ADM
flux derivatives.  No assertion of surjectivity of the constraint map is
involved.

Fix \(0<\beta<\min(q,1)\), and extend \(r\) to a positive smooth
function.  Lemma~\ref{en:strictification} supplies a smooth positive
\(\phi=O_2(r^{-1})\) such that
\begin{equation}\label{va:strict-lapse}
 -\Delta_g\phi=b_*\sqrt{|d\phi|_g^2+r^{-4}}+r^{-3-\beta},
 \qquad \partial_\nu\phi=-1\quad\text{on }B,
 \qquad b_*\ge |K|_g,
\end{equation}
where \(b_*=O(r^{-1-q})\) has smooth radial bounds on the end.
Its Laplacian is integrable and its gradient has a finite flux at infinity.
Use the derivative of the conformal path
\((1+t\phi)^4g,(1+t\phi)^2K\), with both forms fixed, as a fifth
direction.  Denote it by \(V_*\) and put \(\rho_*=r^{-3-\beta}\).
Lemma~\ref{en:conformal} implies that on active rays
\begin{equation}\label{va:strict-cone}
 \mathfrak C'(V_*)\ge8\rho_*,\qquad
 -\theta_\varepsilon'(V_*)=4\quad\text{on }B.
\end{equation}
Moreover \(\mathfrak C'(V_*)/\rho_*\to8\) uniformly over end rays.
The contributions from the drift and the field term are
\(O(r^{-3-q})+O(r^{-5})=o(\rho_*)\).
The four prototypes have the same error bound and hence zero weighted
limit.  Compact directions have zero weighted limit as well.

\begin{lemma}[Constraint separation]\label{va:separation}
There are a probability measure \(\pi_B\) on \(\mathcal M_B\), a
nonnegative locally finite measure \(\Lambda_B\) on \(\mathcal A_B\),
and a nonnegative finite measure \(\zeta_B\) on \(B\), such that
for every compact direction and every ADM prototype,
\begin{equation}\label{va:multiplier}
 16\pi\mathcal L'-c\int_{\mathcal M_B}a'_M(h)\,d\pi_B(M)
 =\int_{\mathcal A_B}\mathfrak C'\,d\Lambda_B
                         -\int_B\theta_\varepsilon'\,d\zeta_B.
\end{equation}
The measure satisfies \(\int\rho_*\,d\Lambda_B<\infty\).
\end{lemma}
\begin{proof}
One-point compactify the active-ray space at infinity.  If it is already
compact, adjoin an isolated point instead.  Assign the limiting values
just calculated to the direction functions.  The boundary and minimizing
family are compact.  Finite linear combinations \(V\) of all the
directions give continuous functions on their disjoint union through
\begin{equation}\label{va:separation-map}
 V\longmapsto\left(
 \frac{\mathfrak C'(V)}{\rho_*},\ -\theta_\varepsilon'(V),\
 \left[c\,a'_M(h)-16\pi\mathcal L'\right]_{M\in\mathcal M_B}
 \right).
\end{equation}
We claim its linear image misses the open cone of functions that are
strictly positive everywhere.

Indeed, a direction in that cone has positive \(V_*\)-coefficient, by
its value at infinity.  Add its compact and prototype changes linearly
and apply the indicated positive conformal change.  The forms stay closed
and retain their end fluxes.  On a compact region, strict positivity on
the closed active set makes the first constraint derivative uniformly
positive there.  Continuity gives this conclusion in a neighborhood of
that set; outside the neighborhood the unperturbed constraint has a
positive minimum.  A uniform Taylor remainder therefore gives the DEC
for all sufficiently small positive parameters.  The same argument gives
the strict chosen boundary expansion.

Here the far end requires a uniform estimate, rather than compact
continuity.  The conformal identities retain a positive multiple of
the original nonnegative ball constraint.  Their additional term is a
positive constant times \(t\rho_*\).  The nonconformal linear errors are
\(tO(r^{-3-q}+r^{-5})\).  The quadratic metric and tensor prototype
errors are \(O(t^2r^{-4})\); the remaining conformal and mixed errors
are \(O(t^2\rho_*)\), after the background nonnegative term has been
retained.  Since \(\beta<q\) and \(\beta<1\), first choose the end radius
large enough for the linear errors to be absorbed, and then choose
\(t>0\) small enough for the quadratic errors to be absorbed.  This
proves feasibility simultaneously on every end ray.  The estimates also
show integrability of the changed constraints and their Taylor
remainders.  The conformal gradient flux and the explicit prototype
fluxes give finite, differentiable ADM limits.  The weak decay,
completeness, and cut hypotheses are preserved.  Thus these paths belong
to the numerical class; no preservation of equality-subclass
outermostness is needed.

Equation~\eqref{va:area-derivative} and positivity of the third component
in Equation~\eqref{va:separation-map} imply
\(c\,a'_+(0)>16\pi\mathcal L'\).  This contradicts the numerical bound
and Equation~\eqref{va:mass-area-derivative}.  The claim follows.

Separation of a linear subspace from a disjoint open convex cone gives
a nonzero continuous functional that annihilates the subspace and is
nonnegative on positive functions.  The representation of positive
functionals on continuous functions gives positive measures on the three
compact test spaces.  If the measure on \(\mathcal M_B\) had zero mass,
evaluation on \(V_*\), using Equation~\eqref{va:strict-cone} also at the
adjoined point, would be strictly positive.  This contradicts
annihilation.  Normalize that mass to one.  On the finite active rays
divide the resulting measure by \(\rho_*\), obtaining \(\Lambda_B\).
The possible atom at infinity contributes zero to compact tests and the
four prototypes.  Rearranging the annihilation identity proves
Equation~\eqref{va:multiplier} and the weighted finiteness assertion.
\end{proof}

\subsection{Normal variations and the minimizing sheets}

The next step uses a Lorentzian metric only as a device for specifying
finite jets.  It does not use existence of a stationary development, or
of an evolution of the additional matter.

\begin{lemma}[Normal test variation]\label{va:gaussian-jets}
For each smooth compactly supported function \(s\) in the open exterior
of \(B\), there are compact variations, indexed by \(\delta>0\), for
which \(h'=2sK\), both flux-form variations are closed, and
\(\mathfrak C'\to0\) uniformly on the compact active-ray support.
Consequently,
\begin{equation}\label{va:averaged-trace}
 \int_{\mathcal M_B}\int_{\partial M}
 s\operatorname{tr}_{T_M}K\,dA_g\,d\pi_B(M)=0.
\end{equation}
\end{lemma}
\begin{proof}
On a neighborhood of the support prescribe a Gaussian metric
\(\mathbf G=-dt^2+g_t\) with \(g_0=g\), \(\partial_tg_t|_0=2K\),
and future normal \(n=\partial_t\).  Prescribe a two-form \(\mathbf F\)
whose initial electric and magnetic fields are the given ones.  The
spatial pullbacks are
\(\mathbf F|_{t=0}=\alpha_2\) and
\((*_{\mathbf G}\mathbf F)|_{t=0}=-\alpha_1\).
The spatial closure equations are precisely \(d\alpha_1=d\alpha_2=0\).
The components of \(d\mathbf F=0\) and
\(d(*_{\mathbf G}\mathbf F)=0\) with one normal slot determine all six
normal derivatives of the two-form from its spatial derivatives and the
already prescribed first metric jets.  Thus both Maxwell equations hold
at the slice.  For the graph variation with velocity \(sn\), Cartan's
formula gives explicitly
\begin{equation}\label{va:lapse-flux-variation}
 \alpha'_1=-d(s\mathcal B^\flat),\qquad
 \alpha'_2=d(s\mathcal E^\flat).
\end{equation}
These are compact closed variations.

Let
\begin{equation}\label{va:maxwell-stress}
 \mathsf S_{ab}=2\left(\mathbf F_{ad}\mathbf F_b{}^d
            -\tfrac14\mathbf G_{ab}\mathbf F_{cd}\mathbf F^{cd}\right),
 \qquad T_\delta=\mathbf{Ein}-\mathsf S.
\end{equation}
The constraint equations give
\(T_\delta(n,n)=\mu_m\) and
\(T_\delta(n,\cdot)|_{T\Omega}=J_m\).
Choose the second normal metric jets so that
\begin{equation}\label{va:regularized-stress}
 (T_\delta)_{ij}=\frac{(J_m)_i(J_m)_j}{\mu_m+\delta}.
\end{equation}
This is a free prescription: the dependence of the spatial Einstein
tensor on \(H=\partial_t^2g_t|_0\) is
\(\tfrac12(H-\operatorname{tr}_g(H)g)\).  Its inverse sends a desired
coefficient \(A\) to \(2A-(\operatorname{tr}_gA)g\) in dimension
three.  Smooth time collars realizing these finite jets can be obtained
by a quadratic Taylor polynomial and a cutoff.  The contracted Bianchi
identity and the Maxwell equations at the slice give
\(\nabla^a(T_\delta)_{ab}=0\) there.  In particular they determine the
normal derivatives of its normal contractions from spatial derivatives;
no further evolution equation is imposed.

We justify a limit used in this calculation.  If smooth \(\mu\ge|J|\)
on a compact set, then
\[
 Q_\delta=\frac{J\otimes J}{\mu+\delta}
 \longrightarrow
 Q=\begin{cases}J\otimes J/\mu,&\mu>0,\\0,&\mu=0\end{cases}
 \quad\text{in }C^1.
\]
At a zero of \(\mu\), nonnegativity and domination give
\(d\mu=dJ=0\).  In a fixed smooth frame,
\[
 |dQ_\delta|\le2|dJ|+|d\mu|,\qquad |Q_\delta|\le\mu.
\]
These bounds are uniform in \(\delta\).  On a sufficiently small
neighborhood of the compact zero set the right sides, including the
first derivatives, are uniformly small.  On its complement the
denominator has a positive lower bound and convergence is smooth.
This proves the claim, including differentiability with derivative zero
on the zero set.

Apply it to Equation~\eqref{va:regularized-stress}.  Where \(\mu_m>0\),
the limiting tensor is \(T=\gamma\otimes\gamma/\mu_m\), where
\(\gamma(n)=\mu_m\) and \(\gamma|_{T\Omega}=J_m\).
The covector \(\gamma\) is causal.  At an active ray with positive
\(\mu_m\), one has
\[
 |v|=1,\qquad J_m=-\mu_m v^\flat,\qquad
 \ell=n+v\text{ null},\qquad \gamma(\ell)=0.
\]
Parallel transport \(\ell\) along any spatial curve.  Its pairing with
the causal covector field \(\gamma\) is nonnegative and vanishes at
the point under consideration.  Its derivative there is therefore zero.
It follows that
\((\nabla_eT)(Z,\ell)=0\) for spatial \(e\) and arbitrary \(Z\).
The divergence identity, in an orthonormal frame, now gives
\begin{equation}\label{va:normal-divergence}
 (\nabla_nT_\delta)(n,\ell)
       =\sum_i(\nabla_{e_i}T_\delta)(e_i,\ell)\longrightarrow0.
\end{equation}
At \(\mu_m=0\), both the limiting tensor and its spatial covariant
derivatives vanish, so the same conclusion holds for every active
\(|v|\le1\).

For completeness the varying arguments in this contraction also matter.
Under the graph variation and metric-square-root transport,
\[
 D_tn=\nabla s,\qquad D_tv=v(s)n,\qquad
 D_t\ell=\nabla s+v(s)n.
\]
The tangential term \(sK^\sharp v\) in the derivative of a graph tangent
vector is canceled by its square-root transport.  Differentiating the
constraint therefore gives
\begin{align*}
 \tfrac12\mathfrak C'_\delta
 &=s(\nabla_nT_\delta)(n,\ell)
       +T_\delta(\nabla s,\ell)
       +T_\delta(n,\nabla s+v(s)n).
\end{align*}
At a positive-density active ray, the last term is
\(J_m(\nabla s)+\mu_m v(s)=0\), and the middle term is
\(\delta J_m(\nabla s)/(\mu_m+\delta)\), which tends uniformly to zero.
At a zero-density ray both terms vanish.  Together with the proved
\(C^1\) convergence and Equation~\eqref{va:normal-divergence}, this
proves uniform convergence on the compact active set.  The variations
have no boundary or ADM contribution.  Since \(\Lambda_B\) is locally
finite, Equation~\eqref{va:multiplier} and dominated convergence give
Equation~\eqref{va:averaged-trace}.
\end{proof}

\begin{lemma}\label{va:trace-support}
For \(\pi_B\)-almost every minimizing set,
\(\operatorname{tr}_{T_M}K=0\) on its sheets in the open exterior of
\(B\).
\end{lemma}
\begin{proof}
There is no cancellation between different tangent planes in
Equation~\eqref{va:averaged-trace}.  Indeed, perimeter submodularity and
the common lower bound \(A\) show that unions and intersections of two
minimizers are minimizers.  Their regularity excludes a transverse
crossing.  Lemma~\ref{ct:compatible-planes} further shows that the plane
incidence set is closed: strict convergence, uniform local density, and
graphical regularity prevent a sheet from disappearing or its tangent
plane from changing at a limiting point.  Thus the average of the
tangent-plane area measures has the form
\(d\eta(x)\,\delta_{T(x)}\), with a single plane over each relevant
spatial point.  Equation~\eqref{va:averaged-trace} reads
\[
 \int s(x)\operatorname{tr}_{T(x)}K\,d\eta(x)=0
\]
for every compact smooth \(s\).  Hence its integrand vanishes
\(\eta\)-almost everywhere.  Fubini's Theorem gives the claim on
\(\pi_B\)-almost every frontier.  On each free smooth sheet it then
holds everywhere by continuity.
\end{proof}

\begin{proposition}[Exclusion of larger marginal obstacles]
\label{va:enlarged-obstacle}
There is no full smooth enlarged obstacle \(B\) of the type described at
the start of this section which contains an additional open subset of
\(\Omega\) outside the original filled obstacle and is marginal with
an arbitrary fixed choice of sign on each component.  For the original
obstacle, the area term in Equation~\eqref{va:multiplier} is exactly
\(a'_S(h)=\tfrac12\int_S\operatorname{tr}_S h\,dA\).
\end{proposition}
\begin{proof}
Every minimizing set for an enlarged obstacle is also a minimizing set
for the original obstacle, because both minima equal \(A\).  Choose one
of the frontiers supplied by Lemma~\ref{va:trace-support}.  Away from
\(S\) this is a free original minimizer, so it is smooth and minimal,
including at contact with an interior component of \(B\).  On the
coincidence set with that component, its second graph derivatives agree
almost everywhere with those of \(B\).  Thus \(H_B=0\) there.  Its
chosen marginality gives \(\operatorname{tr}_{T_M}K=0\), independently
of the chosen sign.  Off the coincidence set the same assertion follows
from Lemma~\ref{va:trace-support}.  Continuity consequently gives
\(H=\operatorname{tr}_{T_M}K=0\) throughout the frontier in
\(\operatorname{int}\Omega\).

At contact with \(S\), the obstacle regularity gives a
\(C^{1,\alpha}\cap W^{2,2}\) graph.  On its coincidence set with \(S\)
the second derivatives agree almost everywhere, and the original
condition \(H_S+\operatorname{tr}_S K=0\) applies.  Off that set the
same equation has just been proved.  The frontier therefore satisfies
the uniformly elliptic expansion graph equation almost everywhere
through the contact patch.  Writing its principal part in divergence
form, local elliptic regularity first gives \(C^{2,\alpha}\) and then
smoothness, since the coefficients and the original data are smooth.
The strong comparison principle for two touching ordered solutions of
this graph equation gives local coincidence with \(S\).  The contact
set is consequently both open and closed in connected \(S\).  If it is
nonempty it is all of \(S\), whose area already equals \(A\); no
additional frontier component remains.

A frontier disjoint from \(S\) would instead be a smooth full enclosing
surface in the interior with \(\theta_+=0\), contrary to the original
outermostness hypothesis.  Thus the chosen frontier must be exactly
\(S\).  Its regular bounded filled side is the original obstacle, so
it cannot contain the additional open set of an enlarged obstacle.
This proves the exclusion.  With \(B=S\), the same reasoning applies
to almost every member selected by \(\pi_S\), showing that their area
functionals all equal \(a'_S\).
\end{proof}

\subsection{Regularity and normalization of the multipliers}

Use the original obstacle from now on, and suppress its subscript.  The
scalar and vector first moments of \(\Lambda\) define measures \(u\)
and \(X\) relative to \(dV_g\).  Positivity and support on active rays
mean, initially as measure statements,
\begin{equation}\label{va:complementarity}
 u\ge |X|_g,\qquad u\mu_m+J_m(X)=0.
\end{equation}
Moment notation is used here because the measures will immediately be
shown to have smooth densities.

\begin{lemma}\label{va:smooth-moments}
The moments have smooth densities on \(\overline\Omega\), with no
additional moment mass on \(S\).  In the interior they satisfy
\begin{equation}\label{va:shift-equation}
 \operatorname{sym}\nabla X=-uK.
\end{equation}
Their full first jets obey a closed linear first-order system with
smooth coefficients.
\end{lemma}
\begin{proof}
For compactly supported interior tensor variations \(k\), the part of
\(\mathfrak C'\) depending on \(k\) integrates to
\[
 2\int\left[u(\tau\operatorname{tr}k-K:k)
                  +X^i\nabla^j(k_{ij}-\operatorname{tr}k\,g_{ij})\right]dV.
\]
There is no interior area term by
Proposition~\ref{va:enlarged-obstacle}.  Integration by parts in
distributions and invertibility of spatial trace reversal give
Equation~\eqref{va:shift-equation}.

For a compact interior metric variation \(h\), the second-order
principal term is the scalar curvature adjoint
\(\nabla^2u-(\Delta u)g\).  Connection variation in the momentum
constraint has at most one derivative of \(h\); after integration by
parts it has at most one derivative of \(u,X\).  The variations of the
fixed flux forms' norms and of the transported ray are algebraic.
All other terms in the metric equation are therefore smooth linear
expressions in \(u,X,\nabla u,\nabla X\).  Taking the trace solves for
\(\Delta u\), since in three dimensions
\(\operatorname{tr}(\nabla^2u-(\Delta u)g)=-2\Delta u\), and hence
solves for the full Hessian of \(u\).

Taking a divergence of Equation~\eqref{va:shift-equation} and using
its trace \(\operatorname{div}X=-u\tau\) gives a vector Laplace
equation for \(X\) with only first derivatives of the moments on the
right.  Together with the traced metric equation this is a linear
elliptic system with diagonal Laplace principal symbol.  Local
distributional elliptic regularity makes the moment measures smooth.
To see that no derivatives remain unaccounted for, differentiate the
symmetrized equation and combine its three index permutations.  For
\(C_{ij}=-2uK_{ij}\) this expresses
\(2\nabla_i\nabla_jX_k\) as
\(\nabla_iC_{jk}+\nabla_jC_{ik}-\nabla_kC_{ij}\) plus curvature
contractions with \(X\).  Thus every second derivative of \(X\), as
well as every second derivative of \(u\), is a smooth linear function
of the first jets.  This is the asserted closed system.

It also proves extension through the boundary.  In a smooth collar,
restrict the system for the first jets to each inward normal segment
from a fixed interior collar slice.  It is a linear ordinary
differential equation with coefficients smooth up to \(S\), so its
solution extends smoothly to \(S\).  Smooth dependence on the
tangential initial point gives a smooth extension of all the moments.
Uniqueness of the normal equations preserves the identities between
these extended jets and derivatives of the extended fields.

It remains possible a priori that \(\Lambda\) has a boundary moment
not represented by these densities.  Take a metric test with
\(h=\nabla h=0\) on \(S\) and arbitrary normal second derivatives of
its tangential block.  In collar coordinates the principal term
\(R'_g=\nabla^i\nabla^jh_{ij}-\Delta\operatorname{tr}h-\langle
\operatorname{Ric},h\rangle\) then has arbitrary boundary value,
through \(-\partial_\nu^2\operatorname{tr}_S h\).  The momentum,
electromagnetic, transported-ray, area, and expansion variations vanish
on \(S\).  After integrating the smooth interior moments by parts,
their boundary terms also vanish because \(h\) has zero first jet.
Equation~\eqref{va:multiplier} consequently annihilates every boundary
test against the scalar moment mass.  That mass is zero.  Its vector
moment is dominated by the scalar mass and is zero as well.
\end{proof}

\begin{lemma}\label{va:normalization}
For any \(q_0\) with \(1/2<q_0<\min(q,1)\),
\begin{equation}\label{va:asymptotic-moments}
 u=b^0+O_2(r^{-q_0}),\qquad
 X^i=b^i+O_2(r^{-q_0}).
\end{equation}
Moreover \(u>0\) in \(\operatorname{int}\Omega\).
\end{lemma}
\begin{proof}
Write \(Y=(u,X)\) in the end coordinates.  Inspection of the preceding
adjoint and its prolongation gives
\begin{equation}\label{va:finite-type-end}
 |\partial^2Y|\le C r^{-1-q_0}|\partial Y|
                         +C r^{-2-q_0}|Y|.
\end{equation}
The derivative coefficients contain \(\partial g,K\); the value
coefficients contain \(\partial^2g,\partial K\), their quadratic
products, and squared electromagnetic fields.  These are controlled by
the original \(O_2/O_1\) hypotheses; no extra derivative bound on the
original data is being used.

Here is the growth argument.  On each radial segment, first bound
\(|Y(r)|\) by its value at a fixed sphere plus
\(\int_R^r|\partial Y(t)|dt\).  Substitution in the integrated form
of Equation~\eqref{va:finite-type-end} and Gronwall's inequality with
the integrable kernel \(t^{-1-q_0}\) give uniformly bounded first
derivatives and \(Y=O(r)\).  The same equation then gives
\(\partial^2Y=O(r^{-1-q_0})\), so \(\partial Y\) has a limit on
each ray with error \(O(r^{-q_0})\).  Two ray values on a sphere can
be joined by an arc of length at most \(\pi r\); their derivative
difference is \(O(r^{-q_0})\).  Thus the limiting derivative matrix
\(L\) is independent of the ray, and
\(Y=Lx+O(r^{1-q_0})\).

The inequality \(u\ge|X|_g\) holds in every direction at infinity.
Applying it on opposite rays first forces the linear part of \(u\)
to vanish, and then forces every linear part of \(X\) to vanish.
Consequently \(\partial Y=O(r^{-q_0})\) and
\(Y=O(r^{1-q_0})\).  Reinsert these bounds in
Equation~\eqref{va:finite-type-end} to obtain
\(\partial^2Y=O(r^{-1-2q_0})\).  Its radial integral from infinity
gives \(\partial Y=O(r^{-2q_0})\).  Since \(2q_0>1\), \(Y\) is
bounded.  One more insertion gives
\(\partial^2Y=O(r^{-2-q_0})\),
\(\partial Y=O(r^{-1-q_0})\), and a common constant limit
\(Y_\infty\) with error \(O(r^{-q_0})\).

Integrate the smooth adjoint identity against the four prototypes,
truncated at radius \(R\), and let \(R\to\infty\).  The limiting
scalar boundary term is \(16\pi u_\infty E'\), and the vector
boundary term is \(16\pi X_\infty^iP'_i\).  The errors involve a
decaying coefficient times a prototype or its derivative, or a
derivative of \(Y\) times a prototype, and their sphere integrals
are \(O(R^{-q_0})\).  The volume integrals converge by the weighted
finiteness in Lemma~\ref{va:separation}.  The prototypes vanish near
\(S\), so Equation~\eqref{va:multiplier} identifies these limits with
\(16\pi(b^0E'+b^iP'_i)\).  Their four independent flux values give
\(Y_\infty=(b^0,b^i)\), proving
Equation~\eqref{va:asymptotic-moments}.

If \(u\) vanished at an interior point, nonnegativity would give
\(du=0\) there and domination would give \(X=0\) and \(\nabla X=0\)
there.  The entire first jet would be zero.  Uniqueness along curves
for the closed first-order system in Lemma~\ref{va:smooth-moments}
would make \(u,X\) identically zero on the connected exterior.  This
contradicts \(u_\infty=E/m>0\).
\end{proof}

\subsection{Stationary reconstruction and electromagnetic potentials}

\begin{proposition}\label{va:stationary-data}
On \(\mathbb R\times\operatorname{int}\Omega\) there are smooth
stationary fields
\begin{equation}\label{va:stationary-metric}
 \mathbf G=-u^2dz^2+g_{ij}(dx^i+X^i dz)(dx^j+X^j dz),
 \qquad \xi=\partial_z,
\end{equation}
and a two-form \(\mathbf F\), with the prescribed orientation, which
induce all four original data \((g,K,\mathcal E,\mathcal B)\) on
each constant-\(z\) slice.  The Killing field is nonzero and future
causal, and tends to the unit future timelike vector
\((E/m,-P/m)\) relative to the original asymptotic normal and frame.
Both Maxwell equations hold.  With \(\mathsf S\) as in
Equation~\eqref{va:maxwell-stress},
\begin{align}
 u(\mathbf{Ein}_{ij}-\mathsf S_{ij})
       &=-(J_m)_{(i}X_{j)},\label{va:spatial-stress}\\
 \mathbf{Ein}-\mathsf S
       &=\frac{\mu_m}{u^2}\,\xi^\flat\otimes\xi^\flat.
       \label{va:null-stress}
\end{align}
The coefficient is nonnegative; wherever it is positive, \(\xi\) is
null.  In particular the region
\(W=-\mathbf G(\xi,\xi)=u^2-|X|_g^2>0\) is electrovacuum.
\end{proposition}

\begin{lemma}[Global Killing potentials]\label{va:potentials}
The spatial one-forms
\begin{equation}\label{va:potential-forms}
 \beta_E=u\mathcal E^\flat+(\mathcal B\times X)^\flat,
 \qquad
 \beta_B=u\mathcal B^\flat+(X\times\mathcal E)^\flat
\end{equation}
are globally exact on \(\operatorname{int}\Omega\).  They are the
pullbacks of \(i_\xi\mathbf F\) and
\(i_\xi(*_{\mathbf G}\mathbf F)\), respectively.  Their potentials
extend smoothly to \(S\).
\end{lemma}
\begin{proof}
Define \(\mathbf F\) uniquely by stationarity and its electric and
magnetic contractions with the future slice normal
\(n=u^{-1}(\partial_z-X)\).  With the given orientation its spatial
parts are \(\alpha_2\) and \(-\alpha_1\) for \(\mathbf F\) and
\(*_{\mathbf G}\mathbf F\).  Contraction with \(\xi=un+X\) gives
Equation~\eqref{va:potential-forms}, including both cross-product signs.

In Equation~\eqref{va:multiplier}, vary \(\alpha_1\) by any smooth
compactly supported closed two-form \(\eta\) in the interior.  The
only contribution is
\(-4\int\beta_E\wedge\eta\), hence it is zero.  Varying
\(\alpha_2\) gives the same assertion for \(\beta_B\).  We explain
why this test class implies exactness on arbitrary topology.  Testing
first with exact \(\eta=d\gamma\) proves \(d\beta=0\).
For any embedded oriented loop choose a trivialized oriented normal
disk bundle and a compact disk two-form of integral one, extended
constantly along the loop and by zero through the tube sides.  It is
a compactly supported closed two-form.  Since \(\beta\) is closed,
its pairing with this form equals the period of \(\beta\) around
the central loop: all parallel loops in the tube have the same period.
Thus every such period is zero.  Smooth one-cycles can be represented
by finite sums of embedded loops in a three-manifold, so every period
is zero.  Path integration defines a global potential.  This argument
uses all compact closed two-form variations, not only the exact ones.

The forms in Equation~\eqref{va:potential-forms} are smooth up to
\(S\).  In a collar, integration in the normal direction extends their
potentials smoothly there; local potentials on overlapping charts
differ by the same constants as in the interior, so the extensions
agree.  Finally, stationarity and Cartan's identity give
\(i_\xi d\mathbf F= -d(i_\xi\mathbf F)=0\) and the corresponding
identity for \(*\mathbf F\).  Their spatial exterior derivatives
vanish because \(d\alpha_1=d\alpha_2=0\).  These two conclusions
give \(d\mathbf F=d(*\mathbf F)=0\) on the product.
\end{proof}

\begin{proof}[Proof of Proposition~\ref{va:stationary-data}]
The slice metric is \(g\), and stationarity of
Equation~\eqref{va:stationary-metric} gives its second fundamental form
\(-u^{-1}\operatorname{sym}\nabla X=K\), by
Equation~\eqref{va:shift-equation}.  The definition of \(\mathbf F\)
recovers the two original fields.  Lemma~\ref{va:potentials} proves
Maxwell closure.  Causality and the end normalization follow from
Equations~\eqref{va:complementarity} and
\eqref{va:asymptotic-moments}; \(\xi\) is nonzero because \(u>0\).

We derive the stress assertion to fix its normalization.  Write
\(\mathsf B(L,M)=L:M-\operatorname{tr}L\operatorname{tr}M\).
In a compact metric variation keep the lapse and shift fixed, and set
\(b=-u^{-1}\operatorname{sym}\nabla X\), so \(b=K\) at the base
point.  Integrating the momentum term by parts gives the gravitational
functional
\[
 \int 2\{u\mu+J(X)\}\,dV
 =\int u\{R-\mathsf B(K,K)+2\mathsf B(K,b)\}\,dV.
\]
Its first variation at \(b=K\) is that of
\(\int u\{R+\mathsf B(b,b)\}\,dV\), since their difference is
\(-\int u\mathsf B(K-b,K-b)dV\), whose first derivative is zero.
The Gauss scalar-curvature identity identifies the latter functional,
up to a divergence, with the spacetime scalar action per unit Killing
time of Equation~\eqref{va:stationary-metric}.  Integrating the
scalar-curvature variation twice by parts gives its coefficient in a
covariant spatial metric variation as \(-u\mathbf{Ein}_{ij}\).
This calculation uses the coframe \(dx^i+X^i dz\): its dual spatial
vectors are tangent to the slices, so the indicated Einstein slots are
exactly the spatial ones.

With a fixed flux form, the vector is its contravariant density divided
by \(\sqrt{\det g}\).  Consequently
\[
 \bigl((|\mathcal E|^2+|\mathcal B|^2)dV\bigr)'
 =\left(\mathcal E_i\mathcal E_j+\mathcal B_i\mathcal B_j
 -\tfrac12(|\mathcal E|^2+|\mathcal B|^2)g_{ij}\right)h^{ij}dV.
\]
The coefficient for twice the lapse times this expression is
\(-u\mathsf S_{ij}\).  Also
\(\langle\mathcal E\times\mathcal B,X\rangle dV\) is independent
of the metric when the two vector densities and \(X\) are fixed.
The transported-ray derivative in Equation~\eqref{va:active-rays}
adds \(-J_{m(i}X_{j)}h^{ij}\).  Including a volume variation costs
nothing because \(u\mu_m+J_m(X)=0\).  The zero interior metric
variation in Equation~\eqref{va:multiplier} therefore yields
Equation~\eqref{va:spatial-stress}.

The normal-normal and normal-spatial components of
\(\mathbf{Ein}-\mathsf S\) are \(\mu_m,J_m\) by Gauss--Codazzi
and the original constraints.  If \(\mu_m=0\), the DEC gives
\(J_m=0\), and Equation~\eqref{va:spatial-stress} makes all
components zero.  If \(\mu_m>0\), complementarity and the two
inequalities \(|J_m|\le\mu_m\), \(|X|\le u\) imply
\[
 |X|=u,\qquad J_m=-\frac{\mu_m}{u}X^\flat.
\]
Since \(\xi^\flat(n)=-u\) and its spatial restriction is
\(X^\flat\), these normal components and
Equation~\eqref{va:spatial-stress} give
Equation~\eqref{va:null-stress}.  This also proves its stated support
property.  In these formulas \(\mu_m,J_m\) are geometric densities;
there is no additional factor of \(8\pi\).
\end{proof}

\subsection{The original boundary and positive surface gravity}

\begin{proposition}\label{va:boundary}
The smooth multipliers on the original boundary satisfy
\begin{equation}\label{va:boundary-identities}
 X=u\nu,\qquad \partial_\nu u+K(X,\nu)=\kappa,
 \qquad \kappa=\frac c2>0.
\end{equation}
Either \(u>0\) everywhere on connected \(S\), or \(u=0\) everywhere
there.  In the first case,
\begin{equation}\label{va:horizon-gradient}
 dW=2\kappa X^\flat\quad\text{on }S.
\end{equation}
In the second case, in original outward normal distance \(s\),
\begin{equation}\label{va:bifurcation-expansion}
 u=\kappa s+O(s^2),\qquad X=O(s^2).
\end{equation}
In both cases \(W>0\) on a deleted collar of \(S\).  The coefficient
\(h_*=\mu_m/u^2\) in Equation~\eqref{va:null-stress} is a smooth
nonnegative function supported in a compact subset of
\(\operatorname{int}\Omega\), and extends by zero through \(S\).
\end{proposition}
\begin{proof}
The integration normal of the exterior at its boundary is \(-\nu\).
An arbitrary compact \(K\)-variation \(k\) therefore leaves the
boundary expression
\[
 -2\int_S\{k(X,\nu)-\operatorname{tr}k\,X_\nu\}\,dA
                         -\int_S\operatorname{tr}_S k\,d\zeta=0.
\]
The mixed normal-tangential components first give \(X_{\mathrm{tan}}=0\),
and the tangential trace gives \(d\zeta=2X_\nu dA\).

Next use a compact scalar test \(s\) and the variation
\((h,k)=(2sg,sK)\), with the flux forms fixed.  The differential
part of \(\mathfrak C'\) is
\(4(-\Delta s+K(\nabla s,v))\), by
Lemma~\ref{en:conformal}.  The expansion derivative at a marginal
boundary is \(2\partial_\nu s\).  After the already proved interior
adjoint cancels the interior terms, Equation~\eqref{va:multiplier}
becomes
\begin{align*}
 -2c\int_Ss\,dA
 &=4\int_S\left[u\partial_\nu s
               -(\partial_\nu u+K(X,\nu))s\right]dA
                         -2\int_S\partial_\nu s\,d\zeta.
\end{align*}
The independent boundary value and normal derivative of \(s\) give
\(X_\nu=u\) and \(\partial_\nu u+K(X,\nu)=c/2\).
This proves Equation~\eqref{va:boundary-identities}.

Let \(L_+\) be the expansion linearization for outward normal motions
of \(S\); its principal part is \(-\Delta_S\).  Apply the multiplier
identity to Lie derivatives of all four original data under a compact
vector field that equals \(s\nu\) on \(S\).  These are legitimate
jets up to the boundary: extend the data smoothly into a collar solely
to take the derivative.  The flux-form Lie derivatives are exact.
At an interior active ray, the transported DEC quantity is a
nonnegative function with value zero, so its Lie derivative vanishes.
There is no boundary constraint atom by
Lemma~\ref{va:smooth-moments}.  The area and expansion derivatives are
\(\int_SH_Ss\,dA\) and \(L_+s\), respectively, while the ADM
derivative is zero.  Since \(d\zeta=2u\,dA\), we obtain
\begin{equation}\label{va:boundary-stability-identity}
 cH_S=2L_+^*u.
\end{equation}
Outer area minimization gives \(H_S\ge0\), by nonnegative outward
normal area variations.  Thus \(u\ge0\) is a supersolution of a
smooth uniformly elliptic operator on connected compact \(S\).
The strong minimum principle, applied locally with the bounded
zeroth-order coefficient, implies either strict positivity everywhere
or identical vanishing.  No positive stability eigenvalue is assumed.

If \(u>0\) on \(S\), then \(W=0\) there.  Tangential derivatives
of \(W\) vanish, and the normal component of
Equation~\eqref{va:shift-equation} gives
\[
 \partial_\nu W
 =2u\partial_\nu u-2\langle\nabla_\nu X,X\rangle
 =2u(\partial_\nu u+uK_{\nu\nu})=2\kappa u>0.
\]
This is Equation~\eqref{va:horizon-gradient} and gives the positive
deleted collar.  If \(u=0\) throughout \(S\), then \(X=0\) there
and all its tangential covariant derivatives vanish.  The mixed and
normal components of Equation~\eqref{va:shift-equation} force its
normal derivatives to vanish as well.  Equation~\eqref{va:boundary-identities}
gives \(\partial_\nu u=\kappa\), proving
Equation~\eqref{va:bifurcation-expansion} and again \(W>0\) for
sufficiently small positive \(s\).

Finally Equation~\eqref{va:asymptotic-moments} gives \(W\to1\) at
infinity.  Equation~\eqref{va:null-stress} makes \(\mu_m=0\) wherever
\(W>0\), so \(h_*\) vanishes on both the boundary collar and a far
end neighborhood.  It is smooth in the remaining compact interior
region because \(u>0\) there.  Extension by zero proves the last claim.
\end{proof}
